\chapter{Kwadratische functies en vergelijkingen}\label{ch:g11:quad}

\omterm{def:g11:quad:quadratic}{Kwadratische functies} zijn de eenvoudigste \omterm{def:g10:functions:function}{functies} na de lineaire, en de
eerste waarvan de \omterm{def:g10:functions:graph}{grafieken} werkelijk gebogen zijn. Dit hoofdstuk bouwt de
volledige gereedschapskist voor hen op: de \omterm{thm:g11:quad:canonical}{toppuntsvorm}, de \omterm{def:g11:quad:discriminant}{discriminant}, het
teken van een kwadratische uitdrukking en de meetkunde van de \omterm{def:g11:quad:parabola}{parabool}.

\section{De toppuntsvorm}

\begin{definition}[Kwadratische functie]\label{def:g11:quad:quadratic}
Een \emph{kwadratische functie}\index{kwadratische functie} is een \omterm{def:g10:functions:function}{functie}
van de vorm
\[
f(x) = ax^2 + bx + c, \qquad a \neq 0,
\]
waarbij $a$, $b$, $c$ \omterm{def:g10:numbers:sets}{reële getallen} zijn. De uitdrukking $ax^2 + bx + c$
heet ook een \emph{kwadratische drieterm}.
\end{definition}

\begin{theorem}[Toppuntsvorm]\label{thm:g11:quad:canonical}
Elke \omterm{def:g11:quad:quadratic}{kwadratische functie} laat zich op precies één manier schrijven als
\[
f(x) = a(x - \alpha)^2 + \beta,
\qquad\text{waarbij } \alpha = -\frac{b}{2a} \text{ en } \beta = f(\alpha).
\]
Dit is de \emph{toppuntsvorm}\index{toppuntsvorm} van $f$.
\end{theorem}

\begin{proof}
We \emph{vullen het kwadraat aan}. Omdat $a \neq 0$, halen we die factor uit
de eerste twee termen:
\[
ax^2 + bx + c = a\left(x^2 + \frac{b}{a}x\right) + c .
\]
Nu is $x^2 + \frac{b}{a}x$ het begin van het kwadraat
$\left(x + \frac{b}{2a}\right)^2 = x^2 + \frac ba x + \frac{b^2}{4a^2}$,
zodat $x^2 + \frac{b}{a}x = \left(x + \frac{b}{2a}\right)^2 -
\frac{b^2}{4a^2}$. Invullen geeft
\[
f(x) = a\left(x + \frac{b}{2a}\right)^2 - \frac{b^2}{4a} + c
= a\bigl(x - \alpha\bigr)^2 + \beta,
\]
met $\alpha = -\frac{b}{2a}$ en $\beta = c - \frac{b^2}{4a}$. Invullen van
$x = \alpha$ maakt de kwadratische term nul, dus is $\beta = f(\alpha)$.
Uniciteit: $a(x-\alpha)^2 + \beta$ \omterm{def:g10:algebra:expand}{uitwerken} en de coëfficiënten van $x$
gelijkstellen dwingt $\alpha = -\frac{b}{2a}$ af, en daarna ligt $\beta$
vast.
\end{proof}

\begin{example}\label{ex:g11:quad:canonical}
Zij $f(x) = 2x^2 - 8x + 3$. Dan is $\alpha = -\frac{-8}{2 \times 2} = 2$ en
$\beta = f(2) = 8 - 16 + 3 = -5$, dus
\[
f(x) = 2(x - 2)^2 - 5 .
\]
Controle door uit te werken: $2(x^2 - 4x + 4) - 5 = 2x^2 - 8x + 8 - 5 =
2x^2 - 8x + 3$.
\end{example}

\begin{proposition}[Extremum en verloop]\label{prop:g11:quad:variations}
Zij $f(x) = a(x-\alpha)^2 + \beta$.
\begin{enumerate}
  \item Is $a > 0$, dan is $f$ \omterm{def:g10:functions:variations}{dalend} op $\intoc{-\infty}{\alpha}$ en
        \omterm{def:g10:functions:variations}{stijgend} op $\intco{\alpha}{+\infty}$: $f$ heeft een
        \emph{\omterm{def:g10:functions:extrema}{minimum}} $\beta$, bereikt in $x = \alpha$.
  \item Is $a < 0$, dan is het verloop omgekeerd: $f$ heeft een
        \emph{\omterm{def:g10:functions:extrema}{maximum}} $\beta$ in $x = \alpha$.
\end{enumerate}
\end{proposition}

\begin{proof}
Stel $a > 0$ en zij $\alpha \leq u < v$. Dan is
\[
f(v) - f(u) = a\bigl((v-\alpha)^2 - (u-\alpha)^2\bigr)
= a\,(v - u)\,\bigl((v-\alpha) + (u-\alpha)\bigr).
\]
Elke factor is positief: $a > 0$, $v - u > 0$, en
$(v-\alpha) + (u-\alpha) > 0$ omdat $v > \alpha$ en $u \geq \alpha$. Dus is
$f(v) > f(u)$: $f$ is \omterm{def:g10:functions:variations}{stijgend} op $\intco{\alpha}{+\infty}$. Voor
$u < v \leq \alpha$ is de laatste factor negatief, dus $f(v) < f(u)$: daar
is $f$ \omterm{def:g10:functions:variations}{dalend}. Ten slotte is $f(x) - f(\alpha) = a(x-\alpha)^2 \geq 0$ voor
alle $x$, zodat $\beta = f(\alpha)$ het \omterm{def:g10:functions:extrema}{minimum} is. Het geval $a < 0$ volgt
door het bovenstaande op $-f$ toe te passen.
\end{proof}

\begin{example}
Een boer heeft $100$ meter omheining om een rechthoekig veld tegen een rechte
muur af te bakenen (langs de muur is geen omheining nodig). Is $x$ de
breedte, dan is de omsloten oppervlakte $A(x) = x(100 - 2x) = -2x^2 + 100x$.
Hier is $\alpha = -\frac{100}{2\times(-2)} = 25$ en
$A(25) = 25 \times 50 = 1250$: de oppervlakte is het grootst bij een breedte
van $25$ m, en bedraagt dan $1250$ m$^2$.
\end{example}

\section{Wortels en ontbinding}

\begin{definition}[Wortel, discriminant]\label{def:g11:quad:discriminant}
Een \emph{wortel}\index{wortel} van de drieterm $ax^2+bx+c$ is een oplossing
van de \omterm{def:g10:algebra:equation}{vergelijking} $ax^2 + bx + c = 0$. De
\emph{discriminant}\index{discriminant} van de drieterm is het getal
\[
\Delta = b^2 - 4ac .
\]
\end{definition}

\begin{theorem}[$ax^2+bx+c=0$ oplossen]\label{thm:g11:quad:roots}
\begin{enumerate}
  \item Is $\Delta > 0$, dan heeft de \omterm{def:g10:algebra:equation}{vergelijking} twee verschillende
        \omterm{def:g11:quad:discriminant}{wortels}
        \[
        x_1 = \frac{-b - \sqrt{\Delta}}{2a},
        \qquad
        x_2 = \frac{-b + \sqrt{\Delta}}{2a},
        \]
        en is $ax^2+bx+c = a(x - x_1)(x - x_2)$.
  \item Is $\Delta = 0$, dan heeft de \omterm{def:g10:algebra:equation}{vergelijking} de enkele \omterm{def:g11:quad:discriminant}{wortel}
        $x_0 = -\frac{b}{2a}$, en is $ax^2+bx+c = a(x - x_0)^2$.
  \item Is $\Delta < 0$, dan heeft de \omterm{def:g10:algebra:equation}{vergelijking} geen reële oplossing, en
        laat de drieterm zich niet met reële coëfficiënten \omterm{def:g10:algebra:expand}{ontbinden}.
\end{enumerate}
\end{theorem}

\begin{proof}
Vertrek van de \omterm{thm:g11:quad:canonical}{toppuntsvorm} uit \cref{thm:g11:quad:canonical}:
\[
ax^2+bx+c
= a\left(x + \frac{b}{2a}\right)^2 + c - \frac{b^2}{4a}
= a\left[\left(x + \frac{b}{2a}\right)^2 - \frac{\Delta}{4a^2}\right],
\]
want $c - \frac{b^2}{4a} = \frac{4ac - b^2}{4a} = -\frac{\Delta}{4a} =
a \times \left(-\frac{\Delta}{4a^2}\right)$.

\emph{1.} Is $\Delta > 0$, dan is $\frac{\Delta}{4a^2} =
\left(\frac{\sqrt\Delta}{2a}\right)^2$ en is de haak een verschil van twee
kwadraten:
\[
ax^2+bx+c = a\left(x + \frac{b}{2a} - \frac{\sqrt\Delta}{2a}\right)
\left(x + \frac{b}{2a} + \frac{\sqrt\Delta}{2a}\right)
= a(x - x_2)(x - x_1).
\]
Een product is nul precies wanneer een van zijn factoren nul is, wat de twee
\omterm{def:g11:quad:discriminant}{wortels} geeft; ze zijn verschillend omdat $\sqrt\Delta \neq 0$.

\emph{2.} Is $\Delta = 0$, dan is de haak $\left(x + \frac{b}{2a}\right)^2$,
en die verdwijnt alleen in $x_0 = -\frac{b}{2a}$.

\emph{3.} Is $\Delta < 0$, dan is $-\frac{\Delta}{4a^2} > 0$, zodat de haak
een kwadraat plus een positief getal is: ze verdwijnt nooit, en de
\omterm{def:g10:algebra:equation}{vergelijking} heeft geen reële oplossing. Een ontbinding $a(x-x_1)(x-x_2)$ zou
\omterm{def:g11:quad:discriminant}{wortels} opleveren, dus bestaat er geen.
\end{proof}

\begin{example}\label{ex:g11:quad:solving}
Los $2x^2 - 8x + 3 = 0$ op. Hier is
$\Delta = (-8)^2 - 4 \times 2 \times 3 = 64 - 24 = 40 > 0$, dus zijn er twee
\omterm{def:g11:quad:discriminant}{wortels}:
\[
x_{1,2} = \frac{8 \pm \sqrt{40}}{4} = \frac{8 \pm 2\sqrt{10}}{4}
= 2 \pm \frac{\sqrt{10}}{2}.
\]
Voor $x^2 + x + 1 = 0$: $\Delta = 1 - 4 = -3 < 0$, geen reële oplossing.
\end{example}

\begin{proposition}[Som en product van de wortels]\label{prop:g11:quad:vieta}
Is $\Delta \geq 0$ en zijn $x_1, x_2$ de \omterm{def:g11:quad:discriminant}{wortels} (gelijk wanneer
$\Delta = 0$), dan geldt
\[
x_1 + x_2 = -\frac{b}{a},
\qquad
x_1 x_2 = \frac{c}{a} .
\]
Omgekeerd zijn twee getallen met som $s$ en product $p$ de oplossingen van
$x^2 - sx + p = 0$.
\end{proposition}

\begin{proof}
Werk de ontbinding van \cref{thm:g11:quad:roots} uit:
\[
a(x - x_1)(x - x_2) = a x^2 - a(x_1 + x_2)x + a\,x_1 x_2 .
\]
De coëfficiënten gelijkstellen aan die van $ax^2 + bx + c$ geeft
$-a(x_1+x_2) = b$ en $a\,x_1x_2 = c$. Voor het omgekeerde: is $u + v = s$ en
$uv = p$, dan is $(x-u)(x-v) = x^2 - sx + p$, dus zijn $u$ en $v$ de \omterm{def:g11:quad:discriminant}{wortels}
ervan.
\end{proof}

\begin{example}
Zoek twee getallen met som $10$ en product $21$. Ze lossen
$x^2 - 10x + 21 = 0$ op; $\Delta = 100 - 84 = 16$, met \omterm{def:g11:quad:discriminant}{wortels}
$\frac{10 \pm 4}{2} = 7$ en $3$. \emph{Snel een \omterm{def:g11:quad:discriminant}{wortel} spotten:} omdat de
coëfficiënten van $x^2 - 5x + 4$ voldoen aan $1 - 5 + 4 = 0$, is $1$ een
\omterm{def:g11:quad:discriminant}{wortel}, en is de andere $\frac{c}{a} = 4$ (hun product moet immers
$\frac ca$ zijn).
\end{example}

\section{Het teken van een drieterm}

\begin{theorem}[Teken van een drieterm]\label{thm:g11:quad:sign}
Zij $f(x) = ax^2+bx+c$ met $a \neq 0$.
\begin{enumerate}
  \item Is $\Delta > 0$, met \omterm{def:g11:quad:discriminant}{wortels} $x_1 < x_2$: dan heeft $f(x)$ het teken
        van $a$ buiten $\intcc{x_1}{x_2}$ en het teken van $-a$ op
        $\intoo{x_1}{x_2}$.
  \item Is $\Delta = 0$: dan heeft $f(x)$ het teken van $a$ voor alle
        $x \neq x_0$, en verdwijnt $f$ in $x_0$.
  \item Is $\Delta < 0$: dan heeft $f(x)$ het teken van $a$ voor alle $x$.
\end{enumerate}
Kortom: \emph{een drieterm heeft het teken van $a$, behalve tussen zijn
\omterm{def:g11:quad:discriminant}{wortels}}.
\end{theorem}

\begin{proof}
\emph{1.} Schrijf $f(x) = a(x-x_1)(x-x_2)$ en volg de tekens van de
factoren. Is $x < x_1$: dan zijn $x - x_1$ en $x - x_2$ allebei negatief, is
hun product positief, en heeft $f(x)$ het teken van $a$. Is
$x_1 < x < x_2$: dan hebben de factoren tegengestelde tekens, is het product
negatief, en heeft $f(x)$ het teken van $-a$. Is $x > x_2$: dan zijn beide
factoren positief en heeft $f(x)$ opnieuw het teken van $a$.

\emph{2.} $f(x) = a(x - x_0)^2$, en $(x-x_0)^2 > 0$ voor $x \neq x_0$.

\emph{3.} Uit het bewijs van \cref{thm:g11:quad:roots} volgt
$f(x) = a\left[\left(x + \frac{b}{2a}\right)^2 +
\left(-\frac{\Delta}{4a^2}\right)\right]$, en de haak is altijd positief.
\end{proof}

\begin{method}[Kwadratische ongelijkheden oplossen]\label{met:g11:quad:inequalities}
Om $ax^2 + bx + c \leq 0$ (of $\geq$, $<$, $>$) op te lossen:
\begin{enumerate}
  \item bereken $\Delta$ en de eventuele \omterm{def:g11:quad:discriminant}{wortels};
  \item stel de \omterm{def:g10:algebra:signtable}{tekentabel} op met \cref{thm:g11:quad:sign} (het teken van $a$
        buiten de \omterm{def:g11:quad:discriminant}{wortels});
  \item lees het oplossingsinterval of de oplossingsintervallen af, met oog
        voor de vraag of de \omterm{def:g11:quad:discriminant}{wortels} zelf meetellen (ruime ongelijkheid) of
        niet (strikte).
\end{enumerate}
\end{method}

\begin{example}\label{ex:g11:quad:inequality}
Los $-x^2 + 3x + 4 \geq 0$ op. Hier is $\Delta = 9 + 16 = 25$, en zijn de
\omterm{def:g11:quad:discriminant}{wortels} $\frac{-3 \pm 5}{-2}$, dus $x_1 = -1$ en $x_2 = 4$. Omdat
$a = -1 < 0$, is de drieterm negatief buiten de \omterm{def:g11:quad:discriminant}{wortels} en positief ertussen.
Bij de ruime ongelijkheid is de oplossingsverzameling $\intcc{-1}{4}$.
\end{example}

\section{De parabool}

\begin{definition}[Parabool]\label{def:g11:quad:parabola}
De \omterm{def:g10:functions:graph}{grafiek} van een \omterm{def:g11:quad:quadratic}{kwadratische functie} $f(x) = a(x-\alpha)^2+\beta$ is een
\emph{parabool}\index{parabool} met \emph{top} $S(\alpha, \beta)$. Ze opent
naar boven wanneer $a > 0$ en naar beneden wanneer $a < 0$, en is
symmetrisch om de verticale rechte $x = \alpha$.
\end{definition}

\begin{remark}
De symmetrie is snel na te gaan: voor elke $h$ is
$f(\alpha + h) = ah^2 + \beta = f(\alpha - h)$, zodat twee punten op gelijke
horizontale afstand van de rechte $x = \alpha$ dezelfde hoogte hebben.
\end{remark}

\begin{omfigure}
\begin{tikzpicture}
\begin{axis}[omaxis,
  width=5.1cm, height=4.8cm,
  xmin=-0.8, xmax=4.4, ymin=-1.7, ymax=3.6,
  xtick={1,3}, xticklabels={$x_1$,$x_2$}, ytick=\empty,
  title={$\Delta > 0$}]
  \addplot[omDef, thick, domain=-0.15:4.15] {(x-1)*(x-3)};
  \fill (axis cs:1,0) circle (1.5pt);
  \fill (axis cs:3,0) circle (1.5pt);
  \fill (axis cs:2,-1) circle (1.5pt) node[below, font=\small] {$S$};
\end{axis}
\end{tikzpicture}%
\begin{tikzpicture}
\begin{axis}[omaxis,
  width=5.1cm, height=4.8cm,
  xmin=-0.8, xmax=4.4, ymin=-1.7, ymax=3.6,
  xtick={2}, xticklabels={$x_0$}, ytick=\empty,
  title={$\Delta = 0$}]
  \addplot[omDef, thick, domain=0.2:3.8] {(x-2)^2};
  \fill (axis cs:2,0) circle (1.5pt)
    node[above right=1pt, font=\small] {$S$};
\end{axis}
\end{tikzpicture}%
\begin{tikzpicture}
\begin{axis}[omaxis,
  width=5.1cm, height=4.8cm,
  xmin=-0.8, xmax=4.4, ymin=-1.7, ymax=3.6,
  xtick=\empty, ytick=\empty,
  title={$\Delta < 0$}]
  \addplot[omDef, thick, domain=0.35:3.65] {(x-2)^2 + 0.8};
  \fill (axis cs:2,0.8) circle (1.5pt) node[below, font=\small] {$S$};
\end{axis}
\end{tikzpicture}

{\small De drie mogelijke liggingen van een naar boven openende \omterm{def:g11:quad:parabola}{parabool}
($a > 0$) ten opzichte van de $x$-as: twee \omterm{def:g11:quad:discriminant}{wortels}, één dubbele \omterm{def:g11:quad:discriminant}{wortel}, geen
\omterm{def:g11:quad:discriminant}{wortel}. De top $S$ heeft \omterm{def:g10:coordgeom:system}{abscis} $\alpha = -\frac{b}{2a}$.}
\end{omfigure}

\begin{omfigure}
\begin{tikzpicture}
\begin{axis}[omaxis,
  width=8.6cm, height=5.6cm,
  xmin=-2.2, xmax=5.4, ymin=-4.4, ymax=6.2,
  xtick={-1,4}, xticklabels={$x_1$,$x_2$}, ytick=\empty]
  \addplot[name path=par, omDef, thick, domain=-1.75:4.75, samples=80]
    {-0.8*(x+1)*(x-4)};
  \path[name path=xaxis] (-1.75,0) -- (4.75,0);
  \addplot[omDef!20] fill between[of=par and xaxis,
    soft clip={domain=-1:4}];
  \fill (axis cs:-1,0) circle (1.5pt);
  \fill (axis cs:4,0) circle (1.5pt);
  \node[font=\small, omDef] at (axis cs:1.5,2.0) {$f > 0$};
  \node[font=\small, omThm] at (axis cs:-1.75,-1.3) {$f < 0$};
  \node[font=\small, omThm] at (axis cs:4.75,-1.3) {$f < 0$};
\end{axis}
\end{tikzpicture}

{\small Het teken van een naar beneden openende drieterm ($a < 0$) met twee
\omterm{def:g11:quad:discriminant}{wortels}: het teken van $-a$ tussen de \omterm{def:g11:quad:discriminant}{wortels} (gearceerd), het teken van $a$
daarbuiten.}
\end{omfigure}

\begin{method}[Een parabool lezen]\label{met:g11:quad:reading}
Om $f(x) = ax^2+bx+c$ te schetsen: bepaal de opening (het teken van $a$), de
top $\left(-\frac{b}{2a},\, f\!\left(-\frac{b}{2a}\right)\right)$, het
snijpunt $(0, c)$ met de $y$-as en de eventuele \omterm{def:g11:quad:discriminant}{wortels}. De symmetrieas
$x = -\frac{b}{2a}$ plaatst daarna elk berekend punt twee keer.
\end{method}

\begin{example}
Voor $f(x) = x^2 - 4x + 3$: opent naar boven, top $(2, -1)$, snijdt de
$y$-as in $(0,3)$, \omterm{def:g11:quad:discriminant}{wortels} $1$ en $3$ (gespot: $1 - 4 + 3 = 0$). Door de
symmetrie ligt ook het punt $(4, 3)$ op de kromme, als spiegelbeeld van
$(0, 3)$.
\end{example}

\section{Oefeningen}

\begin{exercise}[$\star$]\label{exo:g11:quad:1}
Los de \omterm{def:g10:algebra:equation}{vergelijkingen} op:
\[
x^2 - 5x + 6 = 0, \qquad
2x^2 + 3x - 2 = 0, \qquad
x^2 + 2x + 5 = 0, \qquad
9x^2 - 6x + 1 = 0 .
\]
\end{exercise}

\begin{exercise}[$\star$]\label{exo:g11:quad:2}
Schrijf in \omterm{thm:g11:quad:canonical}{toppuntsvorm} en geef de top van de \omterm{def:g11:quad:parabola}{parabool}:
\[
f(x) = x^2 - 6x + 11, \qquad
g(x) = -2x^2 + 4x + 1, \qquad
h(x) = 3x^2 + 6x .
\]
\end{exercise}

\begin{exercise}[$\star$]\label{exo:g11:quad:3}
Los de ongelijkheden op:
\[
x^2 - 3x - 10 < 0, \qquad
2x^2 + x + 3 > 0, \qquad
-x^2 + 6x - 9 \geq 0 .
\]
\end{exercise}

\begin{exercise}[$\star$]\label{exo:g11:quad:4}
Zoek zonder $\Delta$ te berekenen één voor de hand liggende \omterm{def:g11:quad:discriminant}{wortel}, en
daarna de andere:
\[
x^2 - 7x + 6 = 0, \qquad
2x^2 + 5x - 7 = 0, \qquad
x^2 + 4x - 5 = 0 .
\]
(Tip: probeer $x = 1$ en $x = -1$, en gebruik daarna het product van de
\omterm{def:g11:quad:discriminant}{wortels}.)
\end{exercise}

\begin{exercise}[$\star\star$]\label{exo:g11:quad:5}
Zoek twee getallen met som $14$ en product $45$. Zoek daarna de afmetingen
van een rechthoek met omtrek $34$ en oppervlakte $60$.
\end{exercise}

\begin{exercise}[$\star\star$]\label{exo:g11:quad:6}
Zij $f(x) = x^2 + (m-2)x + 1$, met $m$ een reële parameter. Voor welke
waarden van $m$ heeft de \omterm{def:g10:algebra:equation}{vergelijking} $f(x) = 0$ precies één oplossing? En
geen enkele?
\end{exercise}

\begin{exercise}[$\star\star$]\label{exo:g11:quad:7}
Een bal wordt omhoog geworpen; zijn hoogte na $t$ seconden is
$h(t) = -5t^2 + 20t + 1$ (in meter). Wat is de grootste bereikte hoogte, en
op welk tijdstip? Op welk tijdsinterval bevindt de bal zich op minstens
$16$ m hoogte?
\end{exercise}

\begin{exercise}[$\star\star$]\label{exo:g11:quad:8}
Los de \omterm{def:g10:algebra:equation}{vergelijking} $x^4 - 13x^2 + 36 = 0$ op door $X = x^2$ te stellen (een
\emph{bikwadratische} \omterm{def:g10:algebra:equation}{vergelijking}).
\end{exercise}

\begin{exercise}[$\star\star$]\label{exo:g11:quad:9}
Voor welke waarden van $x$ ligt het punt $M(x, x^2)$ van de \omterm{def:g11:quad:parabola}{parabool}
$y = x^2$ strikt onder de rechte $y = x + 2$? Duid het op een schets.
\end{exercise}

\begin{exercise}[$\star\star$]\label{exo:g11:quad:10}
De som van een positief getal en zijn omgekeerde is $\frac{13}{6}$. Zoek het
getal. (Herleid tot een kwadratische \omterm{def:g10:algebra:equation}{vergelijking}.)
\end{exercise}

\begin{exercise}[$\star\star\star$]\label{exo:g11:quad:11}
Zij $\mathcal P$ de \omterm{def:g11:quad:parabola}{parabool} $y = x^2$ en $d_m$ de rechte $y = mx - 1$, met
$m$ een reële parameter.
\begin{enumerate}
  \item Voor welke waarden van $m$ snijden $\mathcal P$ en $d_m$ elkaar in
        twee punten? In één punt? In geen enkel punt?
  \item Bereken, wanneer ze elkaar in precies één punt raken, het raakpunt.
        (In \cref{ch:g11:deriv} zul je $d_m$ herkennen als de raaklijn aan
        $\mathcal P$ in dat punt.)
\end{enumerate}
\end{exercise}

\section{Opgave: het geheime punt van de parabool}

\begin{problem}[{Weekendopgave --- brandpunt en richtlijn: waarom
schotelantennes, koplampen en hangbruggen alle dezelfde kromme tekenen,
waarvan er in wezen maar één bestaat}]
\label{pb:g11:quad:1}
Elke schotelantenne op elk balkon richt haar ontvanger op één welbepaald
punt binnenin; elke koplamp verbergt haar lampje op dezelfde bijzondere
plek. Dat punt --- het \emph{\omterm{pb:g11:quad:1}{brandpunt}} --- is een geheim dat de \omterm{def:g11:quad:parabola}{parabool} al
sinds de Griekse meetkunde bewaart, en de afstandsformule volstaat om het
eruit te halen. Onderweg oefent deze opgave de \omterm{def:g11:quad:discriminant}{discriminant}, werpt ze
projectielen over muren, hangt ze een brug op, en bewijst ze een verbluffend
slotfeit: op een vergroting na is \emph{er maar één \omterm{def:g11:quad:parabola}{parabool} ter wereld}.

\medskip
\noindent\textbf{Deel I --- Vlot met drietermen.}
\begin{enumerate}
  \item Los $x^2 - 5x + 6 = 0$ op, daarna $2x^2 + x - 6 = 0$, en ten slotte
        $x^2 + x + 1 = 0$ (\cref{thm:g11:quad:roots}).
  \item Schrijf $2x^2 - 4x - 6$ in ontbonden vorm en in \omterm{thm:g11:quad:canonical}{toppuntsvorm}
        (\cref{thm:g11:quad:canonical}). Samen met de uitgewerkte vorm heb je
        nu drie kostuums van één \omterm{def:g10:functions:function}{functie}: zeg welk kenmerk elk kostuum in één
        oogopslag toont (\omterm{def:g11:quad:discriminant}{wortels}, top, $y$-afsnede).
  \item Los $-x^2 + 4x - 3 \geq 0$ op met de tekenregel
        (\cref{thm:g11:quad:sign}).
  \item Viète als speurder (\cref{prop:g11:quad:vieta}, in het spoor van
        Babylon in \cref{pb:g10:algebra:1}): zoek twee getallen met som $7$
        en product $12$ door hun \omterm{def:g10:algebra:equation}{vergelijking} op te schrijven. En daarna: één
        \omterm{def:g11:quad:discriminant}{wortel} van $x^2 - 5x + c = 0$ is gelijk aan $2$; zoek $c$ en de
        andere \omterm{def:g11:quad:discriminant}{wortel} zonder ook maar iets op te lossen.
  \item Voor welke waarden van $m$ heeft $x^2 + mx + 9 = 0$ een dubbele
        \omterm{def:g11:quad:discriminant}{wortel}? Geen \omterm{def:g11:quad:discriminant}{wortel}? Twee \omterm{def:g11:quad:discriminant}{wortels}?
\end{enumerate}

\medskip
\noindent\textbf{Deel II --- Het \omterm{pb:g11:quad:1}{brandpunt} onthuld.}
Werk met de \omterm{def:g11:quad:parabola}{parabool} $y = x^2$, het punt $F\left(0, \frac14\right)$ en de
horizontale rechte $\delta$ met \omterm{def:g10:algebra:equation}{vergelijking} $y = -\frac14$.
\begin{enumerate}[resume]
  \item Werk voor een punt $P(x, x^2)$ van de \omterm{def:g11:quad:parabola}{parabool}
        $PF^2 = x^2 + \left(x^2 - \frac14\right)^2$ uit en toon aan dat het
        het volkomen kwadraat $\left(x^2 + \frac14\right)^2$ is. Leid $PF$
        af.
  \item Bereken de afstand van $P$ tot de rechte $\delta$, en besluit:
        \emph{elk punt van de \omterm{def:g11:quad:parabola}{parabool} ligt precies even ver van $F$ als van
        $\delta$}. ($F$ is het \emph{brandpunt}\index{brandpunt}, $\delta$ de
        \emph{richtlijn}\index{richtlijn}.)
  \item Bewijs het omgekeerde: een punt $P(x, y)$ met
        $PF = \operatorname{dist}(P, \delta)$ moet aan $y = x^2$ voldoen.
        (Kwadrateer beide afstanden en vereenvoudig.) De \omterm{def:g11:quad:parabola}{parabool} is dus een
        \emph{meetkundige plaats}, net als de cirkel in
        \cref{pb:g10:coordgeom:1}: gelijke afstand tot een punt en een
        rechte, in plaats van tot één punt.
  \item Pas voor de wijdere \omterm{def:g11:quad:parabola}{parabool} $y = \frac{x^2}{4f}$ (met $f > 0$)
        vraag 6 aan om aan te tonen dat het \omterm{pb:g11:quad:1}{brandpunt} $(0, f)$ is en de
        \omterm{pb:g11:quad:1}{richtlijn} $y = -f$. Toepassing: een schotelantenne is $60$ cm breed
        en $9$ cm diep --- haar rand gaat door $(30, 9)$. Zoek $f$: hoe hoog
        boven de top moet de ontvanger hangen?
  \item Waarom schotels werken: stralen die evenwijdig met de as binnenkomen
        weerkaatsen alle \emph{door het \omterm{pb:g11:quad:1}{brandpunt}} (en een lampje in het
        \omterm{pb:g11:quad:1}{brandpunt} werpt een evenwijdige bundel --- koplampen zijn schotels die
        achterstevoren werken). Leg uit waarom de gelijke afstanden van vraag
        7 dat aannemelijk maken (denk aan de stralen als oprukkende
        golffronten), en zeg welk gereedschap uit \cref{ch:g11:deriv} (zie
        \cref{exo:g11:quad:11}) aannemelijkheid in een bewijs omzet.
\end{enumerate}

\medskip
\noindent\textbf{Deel III --- Parabolen aan het werk.}
\begin{enumerate}[resume]
  \item Een bal volgt de baan $y = x - \frac{x^2}{20}$ ($x$ en $y$ in meter).
        Zoek waar hij neerkomt, de horizontale plaats van zijn hoogste punt,
        en de grootste hoogte (\cref{prop:g11:quad:variations}).
  \item Op dezelfde baan staat een muur in $x = 4$. Bereken de hoogte van de
        bal daar: gaat hij over een muur van $3$ m? En over een van $4$ m?
  \item Een paraboolboog overspant $40$ m met een grootste hoogte van $10$
        m: in gecentreerde \omterm{def:g10:coordgeom:system}{coördinaten} is $y = 10 - \frac{x^2}{40}$. Kan een
        vrachtwagen van $7$ m hoog erdoor op $10$ m van het \omterm{prop:g10:coordgeom:midpoint}{midden}? En op
        $15$ m?
  \item De opbrengst van een theater bij een ticketprijs van $p$ euro is
        $R(p) = p\,(120 - 2p)$ (de vraag zakt als de prijs stijgt). Zoek de
        prijs die de opbrengst maximaal maakt, en dat \omterm{def:g10:functions:extrema}{maximum}.
  \item De hoofdkabels van de Golden Gate zijn tot op grote nauwkeurigheid
        parabolisch (een gelijkmatig verdeelde last op het brugdek levert een
        \omterm{def:g11:quad:parabola}{parabool} op --- een vrij hangende ketting \emph{niet}: haar kromme, de
        kettinglijn, vraagt de exponentiële \omterm{def:g10:functions:function}{functie} van jaar 12). Met een
        overspanning van $1\,280$ m en een doorhang van $143$ m volgt de
        kabel $y = 143\left(\frac{x}{640}\right)^2$ vanaf het \omterm{prop:g10:coordgeom:midpoint}{midden}. Hoe
        hoog ligt de kabel boven zijn laagste punt in $x = 320$ m?
\end{enumerate}

\medskip
\noindent\textbf{Deel IV --- Rechten, raaklijnen en de laatste verrassing.}
\begin{enumerate}[resume]
  \item Vervolledig \cref{exo:g11:quad:11}: de rechten $y = mx - 1$ door het
        uitwendige punt $(0, -1)$ ontmoeten $y = x^2$ naargelang het teken
        van een \omterm{def:g11:quad:discriminant}{discriminant}. Zoek de twee raaklijnen vanuit $(0, -1)$ en hun
        raakpunten.
  \item En nu de rechten $y = mx + 1$ door het \emph{inwendige} punt
        $(0, 1)$: bereken de \omterm{def:g11:quad:discriminant}{discriminant} van de snijvergelijking en toon aan
        dat elke zulke rechte de \omterm{def:g11:quad:parabola}{parabool} twee keer snijdt. Moraal: van
        binnenuit geen raaklijnen --- net als bij een cirkel.
  \item De brandpuntskoorde: snijd $y = x^2$ met de horizontale rechte door
        het \omterm{pb:g11:quad:1}{brandpunt}, $y = \frac14$. Hoe lang is die koorde? Ga de algemene
        regel na (het onthouden waard voor schotels): voor
        $y = \frac{x^2}{4f}$ heeft de koorde door het \omterm{pb:g11:quad:1}{brandpunt} evenwijdig
        met de \omterm{pb:g11:quad:1}{richtlijn} lengte $4f$.
  \item De laatste verrassing: pas de vergroting $(x, y) \mapsto (kx, ky)$
        toe op de \omterm{def:g11:quad:parabola}{parabool} $y = x^2$ en toon aan dat het \omterm{def:g10:functions:function}{beeld} precies
        $y = \frac{x^2}{k}$ is. Besluit: elke \omterm{def:g11:quad:parabola}{parabool} $y = \frac{x^2}{4f}$
        is een vergroting van de eenheidsparabool --- anders dan cirkels en
        ellipsen met verschillende verhoudingen zijn \emph{alle parabolen
        gelijkvormig}.
  \item Slotstuk: het portret van de \omterm{def:g11:quad:parabola}{parabool} in vijf regels --- een
        \omterm{def:g10:algebra:equation}{vergelijking} met drie kostuums (vraag 2), een meetkundige plaats
        (vragen 7--8), de weerkaatser van de ingenieur (vragen 9, 10), de
        baan van een projectiel en de kabel van een brug (vragen 11--15), en
        op een vergroting na één enkele kromme (vraag 19). Telkens één zin.
\end{enumerate}
\end{problem}
